\section{Conclusion}
\label{sec:conclusion}

Algorithmic redistricting does not eliminate gaming opportunities;
it restructures them.
A hostile actor cannot use an algorithm to gerrymander in the traditional
sense---the algorithm cannot see partisan data.
But a hostile actor with control over the deployment context can attempt to
exploit parameter choices, input data, geographic definitions, result selection,
and version selection.

This paper's taxonomy shows that all five gaming vectors are bounded and
detectable under the \dia{} pre-registration and audit protocol.
The plan-affecting vectors (V1--V3, V5) produce a combined maximum shift
of $\leq 0.5$ Democratic seats nationally under simultaneous optimization---a
shift that is (a) smaller than the algorithm's seed-variance noise floor,
(b) far smaller than the 18-seat algorithm-structure effect that B.0 shows
is controlled by statutory pre-specification, and (c) detectable via the
SHA-256 audit chain.

The practical implication for the \dia{}: the pre-registration protocol and
audit chain are not merely bureaucratic safeguards.
They are the mechanism by which algorithmic neutrality is converted from
a property of the computation into a verifiable guarantee about the entire
deployment context.
Without pre-registration, parameter gaming is possible; with it, the
maximum achievable effect is below the noise floor.

Three directions for future work follow from this taxonomy.
R.1 provides a complete adversarial analysis of V1 (parameter gaming)
using U.2's 50-state dataset.
R.2 provides a detailed cryptographic analysis of the V2 (input data)
attack surface, including the TLS limitation and alternative verification paths.
R.4 formalizes Theorem~\ref{thm:gaming-impossibility} as a full
cryptographic argument suitable for Daubert-standard expert testimony,
with worked examples for North Carolina, Texas, and Wisconsin.
